\documentclass{article}
\usepackage[T1]{fontenc}
\usepackage{lmodern}
\usepackage{graphicx}
\usepackage{amsmath,amssymb}
\usepackage[a4paper,margin=2cm]{geometry}
\usepackage[absolute,overlay]{textpos}
\setlength{\TPHorizModule}{1mm}
\setlength{\TPVertModule}{1mm}
\usepackage[indonesian]{babel}
\usepackage{hyperref}
\setlength{\emergencystretch}{2em}
% unit: Halaman 68

\begin{document}

% === Halaman 68 (python/native) ===

•

Playfair cipher termasuk ke dalam kelompok polygram cipher 

•

Pada polygram cipher, satu blok huruf plainteks disubstitusi dengan 

satu blok huruf cipherteks. 

•

Jika satu blok panjangnya 2 huruf, maka ia disebut digram (bigram), 

jika 3 huruf disebut trigram, dst. Playfair cipher melakukan 

enkripsi/dekripsi dalam bentuk bigram.

      Misalnya, bigram  AS diganti dengan RT, BY diganti dengan SL 

68

\end{document}
